% ---
% id: reductions/thm:instrument-instance
% kind: thm
% title: "The branching instance: an instrument followed by branch instances"
% summary: >-
%   A generalized instrument followed by a reduction instance in each branch
%   is a reduction instance, keeping or forgetting the branch label. A branch
%   entry with constant c becomes a zero-constant entry with c times the
%   branch projector moved into the operator, and the normalizer acts branch
%   by branch.
% prominence: standard
% status: accepted
% uses: [def:reduction-instance, def:joint-family, def:exactadm, def:emitnorm, def:generalized-instrument, lem:master, lem:linear, lem:sep, cor:blockdiag, cor:pullback, thm:licence]
% mentions: [cor:register-instance, cor:flag-instance, rem:master-scope]
% assumes: []
% refs: []
% provenance: newly-derived
% lean:
% notes: "The joint-outcome form accumulates over |V| x |J|, so eq:latsumcond must hold at that length wherever the form is emitted (the interface's lattice_exponent row). Step 6 of the proof identifies the normalizer's steps (1)-(4) of def:emitnorm branch by branch, the last resort (3) reading the number of elements and the class weight."
% ---

A receiver described branch by branch, each branch a reduction instance (the
flag instance of Corollary~\ref{cor:flag-instance} for an uncharacterised
branch, for example), is a reduction instance as a whole. The theorem builds
it on Lemma~\ref{lem:master} in two forms: the joint-outcome form keeps the
branch label in the outcome, and the coarse-grained form forgets it. Both
file the branch label as a block index of the target, so they model branches
the adversary may know (Remark~\ref{rem:master-scope}).
Corollary~\ref{cor:register-instance} is the case of a register read by a
POVM.

\begin{theorem}\label{thm:instrument-instance}
Let $J$ be a finite ordered set, let $\{\Ins_j\}_{j\in J}$,
$\Ins_j\colon\Tcl{\Bsys}\to\Tcl{\Bsys'_j}$, be a
\tref{generalized instrument} (Definition~\ref{def:generalized-instrument}),
$\sum_j\tr\circ\Ins_j=\tr$, given independently of any datum, and for each
$j$ let $\mathcal{J}_j$ be a reduction instance on $\Bsys'_j$ over one common
joint alphabet with classes and weights $\big(\Vset,\{\Vset_t\},\{q_t\}\big)$,
with handed datum $A_j$, data sectors $\SecSet_{j,\mathrm{dat}}$ with supports
$\Vset^{(j,s)}$, scalar sectors $\SecSet_{j,\mathrm{sc}}$ with values
$\phi_{j,s}$, hypotheses (H$_j$), channel $\Lambda_j[\cdot]$, transported set
$\hat O^{(j)}[\cdot]$ and entry list $\CList_j[\cdot]$. Write
$\SecSet\coloneqq\bigsqcup_{j\in J}\{j\}\times\SecSet_j$, ordered by $j$ and
then as $\SecSet_j$ is, $\Pi_{(j,s)}$ for its sector projectors,
$\Pi_j\coloneqq\sum_{s\in\SecSet_j}\Pi_{(j,s)}$ for the \emph{branch
projectors}, $\hat X|_j$ for the operator equal to $\hat X$ in block $j$ and
zero elsewhere, and $A\coloneqq(A_j)_{j\in J}$ for the joint datum. Then each
of the two constructions (i) and (ii) is a reduction instance, with
\begin{itemize}[leftmargin=2.2em]
\item the channel
  \begin{equation}\label{eq:instrument-channel}
    \Lambda[A']\;\coloneqq\;\PinS\circ\bigoplus_{j\in J}\Lambda_j[A'_j]\circ\Ins_j ,
  \end{equation}
\item the sector set $\SecSet$, the hypotheses (H)~$=\bigcup_j$(H$_j$)
  together with the hypothesis that $\{\Ins_j\}$ is a generalized
  instrument and, for every $j$ in a declared set $J_{\mathrm{w}}\subseteq J$
  of branches with declared lattice constants $q_j$, the
  \emph{branch-weight hypothesis}
  $\Ins_j^{\dagger}(\idon{\Bsys'_j})=q_j\idon{\Bsys}$,
\item the transported set consisting of the sector projectors and, for every
  $j$ and every $\hat X\in\hat O^{(j)}$, the operator $\hat X|_j$, with
  pull-back $\Ins_j^{\dagger}\big(\Lambda_j^{\dagger}(\hat X)\big)$, in
  particular
  \begin{equation}\label{eq:instrument-transport}
    \Lambda^{\dagger}\big(\Pi_{(j,s)}\big)
    =\Ins_j^{\dagger}\big(\Lambda_j^{\dagger}(\Pi_s)\big),
    \qquad
    \Lambda^{\dagger}\big(\Pi_j\big)=\Ins_j^{\dagger}\big(\idon{\Bsys'_j}\big) ,
  \end{equation}
\item the entry list consisting of, for every branch $j$ and every entry
  $(\hat O_k,\bowtie_k,c_k)\in\CList_j$, the entry
  \begin{equation}\label{eq:instrument-entry}
    \big(\hat O_k|_j-c_k\,\Pi_j,\ \bowtie_k,\ 0\big),
  \end{equation}
  together with the \emph{branch-weight entry} $(\Pi_j,=,q_j)$ for every
  $j\in J_{\mathrm{w}}$, where that set is not empty:
\end{itemize}
\begin{enumerate}[label=(\roman*),leftmargin=3.0em]
\item \emph{(the joint-outcome form)} over the joint alphabet
  $\Vset\times J$ of Definition~\ref{def:joint-family}\,(a) whose
  \tref{outcome alphabet} is $\Yset\times J$, ordered with the branch label
  major and the outcome minor, and whose settings are $\Tset$: writing its
  element $(y,j,t)$ as $(v,j)$ with $v=(y,t)\in\Vset$, its
  \tref{declared order} has
  $(v,j)<(v',j')$ iff $t<t'$, or $t=t'$ and $j<j'$, or $t=t'$, $j=j'$ and
  $y<y'$, so that every class $\Vset_t\times J$ is a contiguous run and the
  restriction of the order to $\Vset\times\{j\}$ \emph{is} the order of
  $\Vset$; its classes are $\Vset_t\times J$ and its weights $q_t$. On it, the model map
  $M_{(v,j)}[A']\coloneqq\Ins_j^{\dagger}\big(M^{(j)}_v[A'_j]\big)$; data
  sectors $(j,s)$, $s\in\SecSet_{j,\mathrm{dat}}$, of dimension $d_{j,s}$ and
  support $\Vset^{(j,s)}\times\{j\}$, with
  $A'^{(j,s)}_{(v,j')}\coloneqq\delta_{jj'}A'^{(s)}_{j,v}$; scalar sectors
  $(j,s)$, $s\in\SecSet_{j,\mathrm{sc}}$, with
  $\phi_{(j,s)}(v,j')\coloneqq\delta_{jj'}\phi_{j,s}(v)$; and the target
  $F_{(v,j)}[A']=F^{(j)}_v[A'_j]\big|_j$, the branch target in block $j$ and
  zero elsewhere;
\item \emph{(the coarse-grained form)} over $\big(\Vset,\{\Vset_t\},\{q_t\}\big)$
  itself: the model map
  $M_v[A']\coloneqq\sum_{j}\Ins_j^{\dagger}\big(M^{(j)}_v[A'_j]\big)$; the same
  data and scalar sectors with supports $\Vset^{(j,s)}$,
  $A'^{(j,s)}_v\coloneqq A'^{(s)}_{j,v}$ and $\phi_{(j,s)}(v)\coloneqq
  \phi_{j,s}(v)$; and the target $F_v[A']=\bigoplus_jF^{(j)}_v[A'_j]$.
\end{enumerate}
In both forms the class-marginal identity is discharged when every branch
discharges its own, and the flag is set in exactly that case; and the
normalizer of Definition~\ref{def:emitnorm} acts on the composite datum
\emph{branch by branch}: its output at the data sectors of branch $j$ is the
output of the normalizer of $\mathcal{J}_j$ at $A_j$, so that the declared
datum is $A'=(A'_j)_j$ with $A'_j$ the branches' declared data and the
declared model is $M[A']$ at those.
\end{theorem}

\begin{proof}
\emph{Idea.} Lemma~\ref{lem:master} combines the branch relations into one;
each part of a reduction instance is then read off branch by branch, and a
branch entry survives the instrument once its constant is moved into the
operator. Throughout, $A'=(A'_j)_j$ denotes a datum of the composite.

\emph{Step 1 (the data).} In the joint form, sector $(j,s)$ has support
$\Vset^{(j,s)}\times\{j\}$, which meets every class $\Vset_t\times J$ because
$\Vset^{(j,s)}$ meets $\Vset_t$; its elements off the support are zero by the
definition of $A'^{(j,s)}$, and its elements on the support are those of
$A'_j$ at sector $s$. Hermiticity and positivity of the elements are those of
$A'_j$; class completeness over $\Vset_t\times J$ reads
$\sum_{(v,j')\in\Vset_t\times J}\delta_{jj'}A'^{(s)}_{j,v}
=\sum_{v\in\Vset_t}A'^{(s)}_{j,v}=q_t\idon{d_{j,s}}$, which is class
completeness of $A'_j$; and the lattice condition is read entry by entry.
In the coarse form the support is $\Vset^{(j,s)}$ and every condition is
literally the branch's. So in both forms a composite datum is an admissible
datum (resp.\ an exactly admissible datum) if and only if every $A'_j$ is:
$\mathcal{D}_{\mathrm{adm}}=\prod_j\mathcal{D}^{(j)}_{\mathrm{adm}}$. The
handed datum $A=(A_j)_j$ is admissible because each $A_j$ is, and its entries
are formed exactly because each branch's are.

\emph{Step 2 ((I1), the model map and stability).} The joint-outcome
alphabet of (i) is a joint alphabet with classes and weights in the sense of
Definition~\ref{def:joint-family}\,(a): its outcome alphabet $\Yset\times J$
is finite, non-empty and ordered, its settings and weights are those of
$\Vset$, and its declared order is that definition's, the outcome varying
fastest and the setting slowest. For every admissible $A'$,
each $M^{(j)}[A'_j]$ is a POVM on $\Bsys'_j$ by (I1) of $\mathcal{J}_j$ under
(H$_j$), which holds by the branch's stability clause. Lemma~\ref{lem:master}\,(i)
with these branch POVMs gives that
$\{\Ins_j^{\dagger}(M^{(j)}_v[A'_j])\}_{(v,j)\in\Vset\times J}$ is a POVM on
$\Bsys$, the joint-form model; coarse-graining $j$ away is
Lemma~\ref{lem:linear}\,(i) for the stochastic matrix that forgets $j$, so
the coarse-form model is a POVM as well. $\Hs{\Bsys}$ is separable by
hypothesis. Stability: (H$_j$) reads no entry of $A'_j$ by the branch's
clause, and the instrument hypothesis and the branch-weight hypotheses read
no datum because the $\Ins_j$ are given independently of any; so (H) reads
no entry of the composite datum,
and the remaining parts, established in Steps 3 to 5 for every admissible $A'$ under
(H), are established as (I1) demands.

\emph{Step 3 ((I2), the target).} The scalar values $\phi_{(j,s)}$ are the
branches' values, \tref[lattice value]{lattice values} in $[0,1]$ formed as (I2) of the branch
requires, with $\sum_{(v,j')}\delta_{jj'}\phi_{j,s}(v)=\sum_v\phi_{j,s}(v)=1$
in the joint form and $\sum_v\phi_{j,s}(v)=1$ in the coarse form. Read
\eqref{eq:instance-target} at the sectors $(j',s)$: in the joint form the
block $(j',s)$ of $F_{(v,j)}[A']$ is $\delta_{jj'}A'^{(s)}_{j,v}$ on a data
sector and $\delta_{jj'}\phi_{j,s}(v)\idon{}$ on a scalar one, so
$F_{(v,j)}[A']$ is the branch target $F^{(j)}_v[A'_j]$ of
\eqref{eq:instance-target} for $\mathcal{J}_j$ in block $j$ and zero
elsewhere, the $F_{v,j}$ of Lemma~\ref{lem:master}\,(ii); in the coarse
form the block $(j,s)$ of $F_v[A']$ is the block $s$ of $F^{(j)}_v[A'_j]$, so
$F_v[A']=\bigoplus_jF^{(j)}_v[A'_j]$, the target of
Lemma~\ref{lem:master}\,(iii). Block diagonality for $\SecSet$ holds by
construction in both forms, and $\dim\Hs{\Qsys}=\sum_j\dim\Hs{\Qsys_j}$ is
finite.

\emph{Step 4 ((I3), the relation).} By (I3) of each branch,
$\big(M^{(j)}[A'_j];O^{(j)}\big)\dsq_{\Lambda_j[A'_j]}
\big(F^{(j)}[A'_j];\hat O^{(j)}\big)$ with $F^{(j)}$ on the
finite-dimensional $\Qsys_j$. Lemma~\ref{lem:master}\,(ii) with these branch
relations gives $\{M_{(v,j)}[A']\}\dsq_{\Lambda_0}\{F_{(v,j)}[A']\}$ for
$\Lambda_0\coloneqq\bigoplus_j\Lambda_j[A'_j]\circ\Ins_j$; its part (v)
transports $\hat X|_j$ to $\Ins_j^{\dagger}(\Lambda_j^{\dagger}(\hat X))$ for
every $\hat X\in\hat O^{(j)}$; its part (iv) gives
$\Lambda_0^{\dagger}(\Pi_j)=\Ins_j^{\dagger}(\idon{\Bsys'_j})$; and
$\Pi_{(j,s)}=\Pi_s|_j$ with $\Pi_s\in\hat O^{(j)}$ by (I3) of the branch, so
$\Lambda_0^{\dagger}(\Pi_{(j,s)})=\Ins_j^{\dagger}(\Lambda_j^{\dagger}(\Pi_s))$,
which is \eqref{eq:instrument-transport}. Its part (iii) gives the
coarse-grained relation with the same channel. Every operator involved is
block diagonal for $\SecSet$ (Step 3; $\hat X|_j$ is block diagonal for
$\SecSet_j$ inside block $j$ by (I3) of the branch), so
Corollary~\ref{cor:blockdiag}\,(i) gives the same relations, transported set
and pull-backs for $\PinS\circ\Lambda_0=\Lambda[A']$, which has the form (I3)
requires; the transported set contains every $\Pi_{(j,s)}$. All of this for
every admissible $A'$ under (H).

\emph{Step 5 ((I4), the entries).} Let $(\hat O_k,\bowtie_k,c_k)\in\CList_j[A'_j]$.
By (I4) of the branch, $\hat O_k$ is a real combination of transported
operators of $\mathcal{J}_j$ of the form \eqref{eq:instance-entry} over
$\SecSet_j$, with coefficients $\alpha_{k,v,s}$ and a projector combination
$Z_k$. Then $\hat O_k|_j$ is the same combination of the operators
$\hat X|_j$, hence transported, and it has the form \eqref{eq:instance-entry}
over $\SecSet$: in the joint form
$\Pi_{(j,s)}F_{(v,j')}[A']\Pi_{(j,s)}=\delta_{jj'}\big(\Pi_sF^{(j)}_v[A'_j]\Pi_s\big)\big|_j$,
so $\hat O_k|_j=\sum_{(v,j')}\sum_{s}\alpha_{k,(v,j'),(j,s)}
\Pi_{(j,s)}F_{(v,j')}\Pi_{(j,s)}+Z_k|_j$ with
$\alpha_{k,(v,j'),(j,s)}\coloneqq\delta_{jj'}\alpha_{k,v,s}$; in the coarse
form $\hat O_k|_j=\sum_v\sum_s\alpha_{k,v,s}\Pi_{(j,s)}F_v\Pi_{(j,s)}+Z_k|_j$
directly. $Z_k|_j$ and $c_k\Pi_j$ are real combinations of sector
projectors of $\SecSet$ with coefficients on the lattice: a coefficient of
$Z_k|_j-c_k\Pi_j$ is $\beta_s-c_k$, a difference of two integer multiples of
$2^{-\kappa}$ and hence one itself, which is all that is asked of it, since
Definition~\ref{def:reduction-instance}\,(I4) puts the band of
\eqref{eq:latinv} on an \emph{entry} and not on a coefficient and this step
uses nothing about the coefficient beyond its being on the lattice. So
$\hat O_k|_j-c_k\Pi_j$ has the form \eqref{eq:instance-entry} with
$Z=Z_k|_j-c_k\Pi_j$, and it is a real combination of transported operators.
Its pull-back, by \eqref{eq:instrument-transport} and linearity, is
\begin{equation*}
  \Lambda^{\dagger}\big(\hat O_k|_j-c_k\Pi_j\big)
  =\Ins_j^{\dagger}\big(\Lambda_j^{\dagger}(\hat O_k)\big)
   -c_k\,\Ins_j^{\dagger}\big(\idon{\Bsys'_j}\big)
  =\Ins_j^{\dagger}\Big(\Lambda_j^{\dagger}(\hat O_k)-c_k\,\idon{\Bsys'_j}\Big).
\end{equation*}
The branch certificate says
$\Lambda_j^{\dagger}(\hat O_k)-c_k\idon{\Bsys'_j}\bowtie_k0$, and
$\Ins_j^{\dagger}$ is a positive map, being the adjoint of a completely
positive one (Lemma~\ref{lem:sep}\,(i) identifies it; positivity is
\eqref{eq:pairing} tested on positive operators): the image of an operator
$\succeq0$ is $\succeq0$, of one $\preceq0$ is $\preceq0$, and of $0$ is $0$.
So the pull-back $\bowtie_k0$, which is \eqref{eq:certificate} for
\eqref{eq:instrument-entry} at the constant $0$, for every admissible $A'$
under (H). The constant moves into the operator because positivity of
$\Ins_j^{\dagger}$ bounds $\Ins_j^{\dagger}(\Lambda_j^{\dagger}(\hat O_k))$
against $c_k\Ins_j^{\dagger}(\idon{})$, the branch weight, and not against
$c_k\idon{\Bsys}$; the entry \eqref{eq:instrument-entry} is the form in
which the branch certificate survives without a hypothesis on that weight.
For $j\in J_{\mathrm{w}}$ the branch-weight hypothesis of (H) makes the
weight a constant, $\Ins_j^{\dagger}(\idon{\Bsys'_j})=q_j\idon{\Bsys}$, and
the branch-weight entry is listed only there; Lemma~\ref{lem:master}\,(iv)
gives
$\Lambda^{\dagger}(\Pi_j)=q_j\idon{\Bsys}$, an equality certificate at
$c=q_j$ for $\Pi_j$, which has the form \eqref{eq:instance-entry} with every
$\alpha=0$; and then also $\Lambda^{\dagger}(\hat O_k|_j)\bowtie_kc_kq_j\idon{\Bsys}$,
the branch entry at the scaled constant, which is
\eqref{eq:instrument-entry} plus $c_k$ times the branch-weight entry.
The class-marginal identity: in the joint form,
$\sum_{(v,j')\in\Vset_t\times J}M_{(v,j')}[A']
=\sum_{j'}\Ins_{j'}^{\dagger}\big(\sum_{v\in\Vset_t}M^{(j')}_v[A'_{j'}]\big)$,
and in the coarse form $\sum_{v\in\Vset_t}M_v[A']$ is the same sum; where
every branch discharges \eqref{eq:instance-classmarginal} this is
$q_t\sum_{j'}\Ins_{j'}^{\dagger}(\idon{\Bsys'_{j'}})=q_t\idon{\Bsys}$ by the
instrument identity $\sum_j\Ins_j^{\dagger}(\idon{\Bsys'_j})=\idon{\Bsys}$
established in the proof of Lemma~\ref{lem:master}, for every admissible
$A'$ under (H). Where some branch does not discharge it, nothing in this
construction establishes the identity, and the flag is not set.

\emph{Step 6 (the normalizer acts per branch).} Definition~\ref{def:emitnorm}
acts per data sector and per class on the elements of the class within the
sector's support. At sector $(j,s)$ of the joint form the support is
$\Vset^{(j,s)}\times\{j\}$ and the class is $\Vset_t\times J$, so the elements
acted on are $\{A^{(s)}_{j,v}:v\in\Vset_t\cap\Vset^{(j,s)}\}$, exactly the
elements the normalizer of $\mathcal{J}_j$ acts on at sector $s$ and class
$t$, with the same residue $q_t\idon{}-\sum_{v\in\Vset_t}A^{(s)}_{j,v}$,
the same exact traces, the same \tref{declared order} (the order of
$\Vset\times J$ declared in (i) restricts on $\Vset\times\{j\}$ to the order
of $\Vset$, the setting varying slowest in both), hence the same
designation order, the same roundings to nearest and mixing weights of the fidelity snap
(Definition~\ref{def:emitnorm}\,(2)), the same components and hosts, and the
same last resort (Definition~\ref{def:emitnorm}\,(3)), which read the
elements of the class, their number, the class weight, the designation order
and the \tref[structural support]{structural supports} alone, all the same
in the two runs under one run configuration and its strict option; every
step of the two runs therefore coincides, and so do their outputs, their
displacements and their moved classes. The
elements \tref{outside the sector support} are touched by neither run. In
the coarse form the support is $\Vset^{(j,s)}$ and the class is $\Vset_t$,
and the identification is immediate. Hence $A'=(A'_j)_j$, and the declared
model is $M[A']$ at the branches' declared data by \eqref{eq:declmodel}.
\end{proof}
